\begin{figure}[t]
\centering
\begin{tikzpicture}
\begin{axis}[
  xbar, bar width=8pt, bar shift=0pt,
  width=0.62\textwidth, height=0.36\textwidth,
  xmin=0, xmax=140, xtick={0,20,...,100},
  xlabel={pipe busy at 20.08 billion iterations/s (\%)},
  symbolic y coords={alu,clmad,issue,lsu,fma},
  ytick={alu,clmad,issue,lsu,fma},
  yticklabels={logic (ALU),\texttt{CLMAD},issue,load/store,FMA},
  y dir=reverse, xmajorgrids,
  nodes near coords, nodes near coords style={font=\scriptsize,anchor=west},
  point meta=explicit symbolic,
]
\addplot[fill=cA,draw=none] coordinates {(91.2,alu) [91.2\% (ceiling 22.0)]};
\addplot[fill=cB,draw=none] coordinates {(73.2,clmad) [73.2\% (ceiling 27.4)]};
\addplot[fill=inkS!45,draw=none] coordinates {(58.7,issue) [58.7\%] (22.1,lsu) [22.1\%] (13.0,fma) [13.0\%]};
\end{axis}
\end{tikzpicture}
\caption{The roofline of Table~\ref{tab:roofline} as a picture: the logic
  pipe is the binding constraint for the table walk's arithmetic on the
  RTX~PRO~6000, with the carry-less unit second; memory and the FMA pipe
  are far from busy. In brackets, the rate at which each unit alone would
  saturate, in billions of iterations per second. Utilisation comes from
  dynamic instruction counts and measured pipe rates, not hardware
  counters (Section~\ref{sec:profiling}).}
\label{fig:roofline}
\end{figure}
